The adaptive extension retains the equal-protection two-fragment protocol and both native complete-object syntaxes, with actual per-source fallback. Calibration evaluated exactly nine policy/SNR cells: NeST QPSK $1/2,2/3$ at 4 dB; NeST QPSK $2/3,3/4,5/6$ and 16-QAM $1/2,2/3$ at 10 dB; and ARPC QPSK $2/3$ at 4 dB and 16-QAM $2/3$ at 10 dB. Targets are the full legal endpoint of each method, not a guarantee that every source sends it. All original 1000 calibration sources and seeds 4101/4102/4103 were used. The explicit prospective nine-cell revision retained the earlier eight-cell scope as history; it allowed 27000 logical objects, 54000 subframes and 108000 header/body decoder attempts.

\begin{table}[t]
\centering
\caption{Frozen low-SNR $N=2048$ policies. DINO-L is the actual calibration mean over 1000 source averages, each containing three noise draws. These are calibration values, not test performance. ARPC denotes ARPC + LDPC--QAM (adapted).}
\label{tab:adaptivepolicies}
\footnotesize\setlength{\tabcolsep}{3.5pt}
\begin{tabular}{lcrr}
\toprule
Method & SNR (dB) & MCS & DINO-L \\
\midrule
NeST-Com & 4 & QPSK $2/3$ & 0.7208 \\
NeST-Com & 10 & 16-QAM $2/3$ & 0.8282 \\
ARPC (adapted) & 4 & QPSK $2/3$ & 0.2930 \\
ARPC (adapted) & 10 & 16-QAM $2/3$ & 0.5604 \\
\bottomrule
\end{tabular}
\end{table}

The completed calibration froze one policy per method/SNR using actual DINO-L after averaging the three noises within each source. Both methods selected QPSK $2/3$ at 4 dB and 16-QAM $2/3$ at 10 dB (Table~\ref{tab:adaptivepolicies}). NeST's selected calibration means are 0.720805 and 0.828199; ARPC's are 0.292956 and 0.560372. Exact ties follow the registered order: lower modulation order, lower nominal code rate, higher target, then lexical policy identifier. No tie or calibration interval triggered additional candidates, and calibration used no bootstrap. The search sets are asymmetric: NeST has two candidates at 4 dB and five at 10 dB, whereas ARPC has one at each. Equal final MCSs therefore do not imply equal search opportunity. These are choices within the nine specified cells, not evidence of a global optimum or estimates of test performance.

The completed frozen-policy test used external\_common100\_v1 at 4/10 dB with seeds 6201/6202/6203: 1200 objects, 2400 subframes and 4800 header/body decoder attempts. Each method/SNR cell correctly delivered all 300 objects; there were no null profiles, rejections or post-output payload mismatches. Thus the adaptive choices restored successful reception in this sample, whereas the separate fixed 64-QAM $5/6$ complete-object reference rejected all objects at these two SNRs. This change also permits actual per-source fallback, so it is not an isolated channel-code ablation or a claim that all received representations are complete m10/e7. With net source capacities of 2288/4848 bits at 4/10 dB, mean source occupancy is 94.90/94.76\% for NeST and 80.34/95.12\% for ARPC; only one of the 100 NeST sources uses m10 at 10 dB, while ARPC uses e4/e5 at 4 dB and e6 at 10 dB.

\begin{table*}[t]
\centering
\caption{Frozen-policy adaptive test at $N=2048$ on external\_common100. Each mean first averages three noises per source. $\Delta$ is NeST minus ARPC; brackets are the saved direct source-paired 95\% intervals. Negative differences favor NeST for downward metrics. Agreement differences are percentage points. The fixed-full 19-dB reference is separate.}
\label{tab:adaptive2048}
\footnotesize\setlength{\tabcolsep}{3pt}
\begin{tabular}{lrrrrrr}
\toprule
 & \multicolumn{3}{c}{4 dB} & \multicolumn{3}{c}{10 dB} \\
\cmidrule(lr){2-4}\cmidrule(lr){5-7}
Metric & NeST & ARPC & $\Delta$ [95\% CI] & NeST & ARPC & $\Delta$ [95\% CI] \\
\midrule
PSNR (dB) $\uparrow$ & 19.256 & 15.477 & 3.779 [3.493, 4.082] & 21.803 & 17.925 & 3.877 [3.688, 4.079] \\
LPIPS $\downarrow$ & 0.1879 & 0.4023 & -0.2144 [-0.2307, -0.1986] & 0.1212 & 0.2517 & -0.1306 [-0.1363, -0.1250] \\
DISTS $\downarrow$ & 0.1585 & 0.2713 & -0.1128 [-0.1246, -0.1013] & 0.1254 & 0.1822 & -0.0568 [-0.0612, -0.0521] \\
DreamSim $\downarrow$ & 0.1152 & 0.3810 & -0.2658 [-0.2933, -0.2386] & 0.0790 & 0.1965 & -0.1174 [-0.1309, -0.1049] \\
DINO-L $\uparrow$ & 0.7014 & 0.2715 & 0.4298 [0.3885, 0.4718] & 0.8129 & 0.5245 & 0.2884 [0.2474, 0.3304] \\
Agreement (\%) $\uparrow$ & 75.0 & 36.0 & 39.0 [28.0, 50.0] & 85.0 & 73.0 & 12.0 [2.0, 23.0] \\
MS-SSIM $\uparrow$ & 0.6822 & 0.4169 & 0.2653 [0.2456, 0.2854] & 0.8071 & 0.6176 & 0.1895 [0.1789, 0.2001] \\
\bottomrule
\end{tabular}
\end{table*}

Table~\ref{tab:adaptive2048} reports all seven test metrics. NeST-Com's PSNR is 19.256 versus ARPC's 15.477 dB at 4 dB and 21.803 versus 17.925 dB at 10 dB. The direct source-paired gains are 3.779 [3.493, 4.082] and 3.877 [3.688, 4.079] dB. DreamSim reductions are 0.26583 [0.23861, 0.29329] and 0.11740 [0.10491, 0.13090]; agreement gains are 39 [28, 50] and 12 [2, 23] percentage points. All fourteen pointwise paired intervals favor NeST-Com in the appropriate metric direction. They describe these two frozen adapted systems on the matched sources, not a multiplicity-adjusted guarantee or an endpoint-only causal effect.

The test first averaged the three noise draws within each source, then generated exactly one 10000-replicate source-paired bootstrap bank with seed 2026101902. Exact input/pixel/metric identities permitted reuse of 159 of the 400 distinct test image pairs; 241 pairs required new recovery and scoring (199 NeST and 42 ARPC). All 1200 logical frames remain in the denominator. Calibration and test together used the fixed 28200-object, 56400-subframe and 112800-decoder-attempt budget, with no further candidates. The 19-dB fixed-full reference remains separate: its complete-representation results include ARPC's lower DISTS and cross-zero DINO/agreement intervals. The low-SNR adaptive result therefore does not erase the high-SNR budget boundary or imply one calibrated policy across all three SNRs.
